\section{A section on the original normal variety}\label{sec:conclusion}

The geometric argument produces a pluricanonical section on a smooth
fourfold. Hashizume's reduction gives an effective real representative
of an lc adjoint on the original variety. The following elementary
lemma turns that representative into a rational one. It concerns
finite linear algebra and uses no nonvanishing or abundance theorem.

\begin{lemma}[Finite rational feasibility]\label[lemma]{lem:rational-feasibility}
Let $M=(m_{ij})\in\operatorname{Mat}_{w\times v}(\Q)$ and
$d\in\Q^w$. If the system
\begin{equation}\label{eq:rational-feasibility}
 d_i-y_i=\sum_{j=1}^{v}m_{ij}b_j\quad(1\le i\le w),
 \qquad y_i\ge0
\end{equation}
has a real solution, it has a rational solution. Moreover, one may
require precisely the same $y_i$ to be zero as in the given solution.
Consequently, if $D$ is a rational Weil divisor on a normal integral
variety $X$ over a field $k$ and $D\sim_{\R}G$ for an effective real Weil divisor $G$,
then $D\sim_{\Q}G_{\Q}$ for an effective rational divisor with
$\Supp G_{\Q}=\Supp G$. If $rD$ is Cartier for a prescribed
positive integer $r$, then $H^0(X,\cO_X(mrD))\ne0$ for some $m>0$.
\end{lemma}

\begin{proof}
Fix the equations $y_i=0$ at all coordinates vanishing in the given
real solution. Together with the displayed equations they define a
nonempty affine space over $\Q$. Gaussian elimination over $\Q$
gives a rational particular solution and a basis over $\Q$ for its
direction space. Thus its rational points are dense in its real
points. Approximate the given point closely enough to keep all of
the finitely many remaining coordinates $y_i$ strictly positive.
This proves the first assertion, including the case where all
$y_i$ vanish or the affine space has dimension zero.

For the divisor assertion, choose a finite expression
\[
 D-G=\sum_{j=1}^{v} a_j\Div(f_j),\qquad
 a_j\in\R,\quad f_j\in k(X)^*.
\]
List the prime divisors in the supports of
$D,G$, and these finitely many principal divisors as
$P_1,\ldots,P_w$. Write $D=\sum d_iP_i$ and let $m_{ij}\in\Z$
be the coefficient of $P_i$ in $\Div(f_j)$. The coefficients of
$G$ and the $a_j$ give a real solution of
\eqref{eq:rational-feasibility}. A rational solution gives
$G_{\Q}=\sum y_iP_i\ge0$ with the same support and
$D-G_{\Q}=\sum b_j\Div(f_j)$.

Choose a positive integer $N$ divisible by $r$ and by the
denominators of all $y_i,b_j$. Then
\[
 ND-NG_{\Q}=\Div\!\left(\prod_j f_j^{Nb_j}\right).
\]
The effective integral divisor $NG_{\Q}$ is Cartier, since $ND$
is Cartier and their difference is principal. It defines a nonzero
section of $\cO_X(ND)$. Write $N=mr$. No $\Q$-factoriality is
needed: all equalities were equalities of Weil divisors on $X$.
\end{proof}

\begin{proof}[Proof of \Cref{cor:lc-nonvanishing}]
Hashizume's reduction \cite[Theorem~1.4]{Hashizume2018} takes smooth
canonical nonvanishing in dimension four, supplied by
\Cref{thm:nonvanishing}, and gives
nonvanishing for projective lc pairs with effective real boundaries
and pseudo-effective real Cartier adjoints in dimensions at most four.
Applied to the given rational lc pair, it produces $D\sim_{\R}G\ge0$
on the original normal variety; nefness implies pseudo-effectivity.
\Cref{lem:rational-feasibility} gives a nonzero section of $mrD$
for some positive integer $m$, with the prescribed index $r$.
\end{proof}
